\documentclass[10pt]{article}

    

\def\Type{LessonCaelan}
\def\TypeNum{Lesson 0}
\def\TypeTitle{Modeling Change}
\def\Semester{Fall 2021} %Not Used 
\def\Course{MATH 1210}
\def\Targets{\bf Intro to the course!}
\def\desmosClassCode{{\bf ZAJ$3$AZ}}

\usepackage{preambleDJ} 

\begin{document}
\pagestyle{otherpages}
%\thispagestyle{empty}
\thispagestyle{firstpage}


\vs{.75}

{\bf\color{blue}About this class}

\vfill

{\bf\color{blue}Introduction}

A \textit{function} is a deterministic rule for transforming one sort of thing into another. The collection of things which can go into the function is called the \textit{domain} of the function, and the collection of things which can come out is called the \textit{range}.

The \textit{real numbers} comprise everything on the number line, both positive and negative.

\begin{enumerate}
    \item (In groups) Give some examples of functions.\\[0.5cm]
    \begin{tabular}{c | c | c}
         \parbox{0.3\textwidth}{
            Domain: Real numbers\\[7pt]
            Range: Real numbers
            \vspace{4cm}
         }
         &
         \parbox{0.3\textwidth}{
            Domain: \underline{\hspace{3cm}}\\[7pt]
            Range: Real numbers
            \vspace{4cm}
         }
         &
         \parbox{0.3\textwidth}{
            Domain: \underline{\hspace{3cm}}\\[7pt]
            Range: \underline{\hspace{3cm}}
            \vspace{4cm}
         }
    \end{tabular}
\end{enumerate}
\vspace{0.5cm}

\newpage

{\bf What is Calculus?}
\vspace{8cm}

{\bf Calculus and Geometry}

\vspace{6.5cm}



{\bf\color{blue} Application: Environmental Science}\\
\mpageT{.4}{.6}[\hs{.1}]{
The Keeling Curve is the record of atmospheric CO$_2$ in parts per million (ppm). This data set is named after Charles Keeling who started collecting this data at Mauna Loa Observatory, starting in 1958.  A sample of atmospheric CO$_2$ is shown below in five-year intervals from starting in 1965 through 2010, each measurement taken on January 15th of that year.  Note that ``year" means ``years since 01/15/1965."  For example, $t=10$ represents the data collected on 01/15/1975.
}{
\graphicsC{keelingDataMMA}{3.5}
}

\begin{center}
\begin{tabular}{|c|c|c|c|c|c|c|c|c|c|c|}
\hline
year & 0 & 5 & 10 & 15 & 20 & 25 & 30 & 35 & 40 & 45 \\
\hline
CO$_2$ level & 320 & 325 & 331 & 338 & 345 & 354 & 360 & 369 & 379 & 389 \\
\hline
\end{tabular}
\end{center}

\newpage

\enum{
\setcounter{enumi}{1}
\item (In groups)  What's your best guess for the concentration of CO$_2$ in the atmosphere on 01/15/2025?
\vspace{3cm}
\item (In groups) What about 01/15/2030?
\vspace{3cm}
\item (In groups) On average, how much did the concentration of CO$_2$ change each year between 1965 and 2000?

\vfill

\item (In groups) Now, let's take a more detailed look at the Keeling data. You can access it directly at \url{https://scrippsco2.ucsd.edu/data/atmospheric_co2/primary_mlo_co2_record.html} from the Scripps Institution of Oceanography. You can access a copy of the monthly data that has already been loaded into Desmos at \url{https://www.desmos.com/calculator/o2knft7vce}.

How did your guess for Question 2 compare to the true answer?
\vfill
\newpage
\item (In groups) How would you describe atmospheric CO$_2$ concentration as a function of time?

Try using Desmos to fit various functions to the data.
\vfill
\item (In groups) Should you amend your answer to Question 3?
\vfill
\item (In groups) What 10-year period has had the highest average rate of change of CO$_2$ concentration?
\vfill
\item (In groups) Choose a member of your group. Estimate how atmospheric CO$_2$ levels changed over the 24 hours of their most recent birthday.
\vfill
}
% \newpage



% \item  (In groups) Work together to brainstorm at least two strategies to refine your estimate and/or increase your confidence in your estimate. You might find it helpful to use Desmos widget ``First Day of Class: The Keeling Curve"  to implement your ideas. This can be found on the ToDo tab at \url{https://student.desmos.com} using code \desmosClassCode (You might need to create a free account at Amplify). 
% \vs{3}
% \item (In groups) In what ways would you describe how the amount of CO$_2$ is changing each year?\newpage
% \item (Extra, if time. In groups) What other information might be helpful to further  improve your estimate and/or build confidence in your estimate? What mathematical tools might be helpful?
% }
% \vs{3}
% \bblue{Key Questions for the semester}
%\vs{1} \\
%\pe Based on this data, what do you think is a reasonable guess for the amount of CO$_2$ in the atmosphere in January 2024?

\end{document}